\section{Polynomial expected depth with improved initial batches}
\label{sec:first-phase}

The early geometric batches need not retain their original budgets.
Allowing a different small first budget, and sometimes merging the first
two batches, improves the general polynomial expected-depth factor to
below $5.197569$. The later batches and the rational adaptive block modes
of Section~\ref{sec:adaptive-blocks} remain unchanged. In particular, the
improvement applies to arbitrary strictly positive rational priors.

The proof can be read in four steps. First specify the enlarged
polynomial family of actual schedules. Next compute its signed change
to the first-hit and all-negative costs. Then use one finite table to
certify the positive-rank coefficients and every possible zero waiting
length. Finally invoke the same capped-rank and stage-to-tree lemmas
as for the earlier block variants. The table is a comparison
distribution used to bound the algorithm's minimum rational score.

\subsection{A larger polynomial family of schedules}

Use the same full-hull edge candidates, exact budgeted coverage sets
$Q_k$, minimum cover $Q_*$, and greedy ordering within each batch.
Write $k_{\min}$ for the minimum cover size, reserving
$k_j=\lfloor\beta^ju\rfloor$ for a shift's geometric budgets. Set
\begin{equation}\label{eq:first-phase-parameters}
 \beta=\frac{3512863}{10^6},\qquad
 \delta=\frac{155223384}{10^9}.
\end{equation}
For every shift representative $u\in[1,\beta)$, try these sequences:
\[
 \begin{array}{ll}
 q,k_1,k_2,\ldots,&q\in\{1,\ldots,6\},\quad q<k_1,\\
 k_1,k_2,\ldots.&
 \end{array}
\]
The second sequence merges the old first two batches. At \emph{every}
phase, including the first, use $Q_k$ when its budget $k<k_{\min}$;
otherwise use $Q_*$ and finish. Thus an initial budget reaching
$k_{\min}$ starts with the minimum cover and has no later phase.
Omit repeated and unproductive blocks. The original first budget is
included because $k_0\le3$ and $k_0<k_1$.

For each sequence use Algorithm~\ref{alg:block-mode}, with the same
$p_*,h_*$ and rational saving $D$ as in
\eqref{eq:rational-block-saving}. Choose the schedule minimizing
\eqref{eq:adaptive-shift-objective}. The physical stage bound
\eqref{eq:adaptive-stage}, the coarsened transcript
$J_*$, and the individual-cap bound
\eqref{eq:coarse-discovery-cap} therefore still hold.
Algorithm~\ref{alg:first-phase-family} records the new family.

Refine the geometric shift partition of
Algorithm~\ref{alg:batch-orders} by adding every breakpoint
\begin{equation}\label{eq:initial-shift-refinement}
 \frac{k}{\beta},\ \frac{k}{\beta^2}\in[1,\beta),
 \qquad 1\le k\le\lceil\beta^3\rceil.
\end{equation}
Use each breakpoint below $\beta$ and choose one interior point from
each consecutive pair as a representative. This constant
refinement, for fixed $\beta$, makes both $k_1$ and $k_2$ constant
on every initial cell even when the minimum-cover size truncates the
base partition. It also supports the two-budget family in the next
section. All later executed budgets have their breakpoints in the
base partition. Thus each breakpoint and open cell has a representative,
and every schedule and mode score is fixed within each open cell. There
are at most seven variants per representative. The
certificate distributions are used only in the proof.

\begin{algorithm}[H]
\caption{Improved first-phase family for \ExpectedTree}
\label{alg:first-phase-family}
\KwIn{Weighted positive-pattern block demands, zero weight $p_0$;
$\beta,\delta$ from \eqref{eq:first-phase-parameters}.}
\KwOut{A minimum-score supplied-block schedule.}
Reuse the coverage sets $Q_k$ and minimum cover $Q_*$ of size $k_{\min}$\;
Refine the \BatchOrders\ shift partition by
\eqref{eq:initial-shift-refinement}; use its breakpoints and cell midpoints\;
\ForEach{shift representative $u$}{
 $k_1\gets\lfloor\beta u\rfloor$\;
 $\mathcal V\gets\{(q,1):1\le q\le6,\ q<k_1\}\cup\{(k_1,2)\}$\;
 \ForEach{$(q,j_0)\in\mathcal V$}{
  $G\gets$ empty order\;
  Process the first budget $q$: use $Q_q$ if $q<k_{\min}$,
  and $Q_*$ otherwise\;
  Greedily append productive unused blocks by newly covered mass\;
  \If{$q<k_{\min}$}{
   \For{$j=j_0,j_0+1,\ldots$}{
    $k\gets\lfloor\beta^ju\rfloor$;
    process $Q_k$ if $k<k_{\min}$, and $Q_*$ otherwise\;
    Greedily append productive unused blocks\;
    \If{$k\ge k_{\min}$}{stop these phases\;}
   }
  }
  Use Algorithm~\ref{alg:block-mode} to compute rational savings $D$\;
  Retain $G$ with score $2\MSSC_{\rm blk}(G)+2p_0|G|-\sum D$\;
 }
}
\KwRet a retained schedule of minimum score\;
\end{algorithm}

Use this routine in place of \BatchOrders\ in the adaptive variant
of \ExpectedTree, with the new threshold $\delta$ at every recursive
state. All positive exits still learn the full pattern before child
attachment. An empty demand family needs no schedule; the direct
degenerate-hull and singleton cases are handled by \ExpectedTree.

\subsection{Signed first-hit corrections and zero waiting}

The unaveraged accounting bound in the proof of
Lemma~\ref{lem:mssc} is
\[
 \phi_u(t)=\frac{k_0+1}{2}
  +\frac12\sum_{j\ge0}(k_j+k_{j+1})\mathbf1_{\{t>k_j\}}.
\]
For the sequence beginning at $q<k_1$, its counterpart is
\[
 \phi_{u,q}(t)=\frac{q+1}{2}
  +\frac{q+k_1}{2}\mathbf1_{\{t>q\}}
  +\frac12\sum_{j\ge1}(k_j+k_{j+1})\mathbf1_{\{t>k_j\}}.
\]
To compare the schedules, double both invoices. Their common tail
from phases $j\ge1$ cancels, leaving
\begin{equation}\label{eq:first-phase-correction}
 \begin{aligned}
 d_q(u,t)&=2\bigl(\phi_u(t)-\phi_{u,q}(t)\bigr)\\
 &=k_0-q+(k_0+k_1)\mathbf1_{\{t>k_0\}}
              -(q+k_1)\mathbf1_{\{t>q\}}.
 \end{aligned}
\end{equation}
For the merged sequence beginning at $k_1$, its doubled invoice is
\[
 2\phi_{u,\mathrm m}(t)=k_1+1+(k_1+k_2)\mathbf1_{\{t>k_1\}}
       +\sum_{j\ge2}(k_j+k_{j+1})\mathbf1_{\{t>k_j\}}.
\]
Subtracting this from $2\phi_u(t)$ cancels the common tail
$j\ge2$; the remaining terms give
\begin{equation}\label{eq:merged-phase-correction}
 d_{\rm m}(u,t)=k_0-k_1+(k_0+k_1)\mathbf1_{\{t>k_0\}}.
\end{equation}
These corrections can be negative. Retaining their signs accounts
for both accelerated and delayed comparison ranks. The within-batch
greedy proof applies to either sequence, and each resulting first-hit
invoice is nondecreasing. Thus a comparison edge hitting a class
before its supporting-prefix index can only improve its bound.
Early completion at $Q_*$ preserves the formulas as upper bounds. The
unrolled residual expression adds only nonnegative terms after the last
actual phase.

The all-negative route requires a separate correction. The old
schedule uses the minimum cover of size $k=k_{\min}$, together with
each earlier budget below $k$, so its length invoice is
\[
 k_{\min}+\sum_{j:k_j<k_{\min}}k_j,
\]
whereas the regular variant uses
\[
 k+q\mathbf1_{\{q<k\}}+\sum_{j\ge1:k_j<k}k_j.
\]
Subtracting the old invoice cancels the final cover and the common
tail $j\ge1$, leaving $z_q=q\mathbf1_{\{q<k\}}-k_0\mathbf1_{\{k_0<k\}}$.
For the merged variant the invoice is
\[
 k+k_1\mathbf1_{\{k_1<k\}}+\sum_{j\ge2:k_j<k}k_j,
\]
so the common final-cover and $j\ge2$ terms cancel and leave
$z_{\rm m}=-k_0\mathbf1_{\{k_0<k\}}$. Thus the two variants
change the old length accounting by, respectively,
\begin{equation}\label{eq:first-phase-zero-correction}
 \begin{aligned}
 z_q(u,k)&=q\mathbf1_{\{q<k\}}-k_0\mathbf1_{\{k_0<k\}},\\
 z_{\rm m}(u,k)&=-k_0\mathbf1_{\{k_0<k\}},
 \qquad k=k_{\min}.
 \end{aligned}
\end{equation}
In particular, this counts all executed prefix blocks before a
zero-pattern exit, including a first batch that does not complete
the cover.

\subsection{An eight-row rational certificate}

Keep the witness masses $m_t$, their prefix masses
$M_q=\sum_{t\le q}m_t$, and the nondecreasing discounts $e_\beta(t)$
from Section~\ref{sec:polynomial-blocks}. They obey
\[
 \sum_t m_t=b,\qquad m_t\le2\delta w,
 \qquad w=w(S),\quad b=w-p_0.
\]
The first budget-$q$ batch covers at least $M_q$, by exact coverage,
and uses at most $q$ blocks. Lemma~\ref{lem:affine-block-credit} gives
the following minorants for the \emph{total} saving:
\[
 \mathsf A_q=M_q-\delta wq,\qquad
 \mathsf C_q=c_*M_q-d_*\delta wq,\qquad 0.
\]
They hold also when that batch completes $Q_*$ early. As before,
convex mixtures of these bounds count the actual saving once.

Fix
\[
 \begin{aligned}
 v_1&=\frac{95240701}{10^8},&v_2&=\frac{37279986}{10^8},&
 v_3&=\frac{4271137}{10^8},\\
 v_4&=\frac{33933837}{10^8},&v_5&=\frac{36583316}{10^8}.
 \end{aligned}
\]
Average over the log-uniform shift density $1/(u\ln\beta)$.
The table specifies the conditional distribution of variants and
saving minorants. A pair $(q,\mathsf L)$ means first budget $q$
and minorant $\mathsf L$; except for the two marked merged rows,
the remaining budgets are $k_1,k_2,\ldots$.
\begin{center}
\small
\begin{tabular}{ll}
\hline
Shift range & Variant and saving minorant\\
\hline
$[1,4/\beta)$ & $(3,\mathsf C_3)$, merged\\
$[4/\beta,5/\beta)$ & $v_1(4,\mathsf A_4)+(1-v_1)(4,0)$, merged\\
$[5/\beta,2)$ & $(1,\mathsf C_1)$\\
$[2,8/\beta)$ & $v_2(1,\mathsf C_1)+(1-v_2)(2,\mathsf C_2)$\\
$[8/\beta,3)$ & $(2,\mathsf C_2)$\\
$[3,11/\beta)$ & $v_3(2,\mathsf C_2)+(1-v_3)(3,\mathsf C_3)$\\
$[11/\beta,12/\beta)$ & $(3,\mathsf C_3)$\\
$[12/\beta,\beta)$ &
 \shortstack[l]{$v_4(3,\mathsf C_3)+v_5(5,\mathsf A_5)$\\
 $\quad +(1-v_4-v_5)(6,0)$}\\
\hline
\end{tabular}
\end{center}
All conditional weights are nonnegative and sum to one. In the
merged rows $k_1$ is exactly three or four. Every table row whose
$k_1$ varies uses its original first budget $q=k_0$, so its phase
correction is identically zero.

Here are explicit definitions of the analytical coefficients. Begin
with $a_t=e_\beta(t)$ and $\mu=0$. A row $[l,r)$ has probability
$P=\ln(r/l)/\ln\beta$. An option of conditional weight $v$, initial
budget $q$, and minorant $sM_q-\rho\delta w$ contributes
\begin{equation}\label{eq:first-phase-coefficients}
 a_t\gets a_t+Pv\bigl[d_q(u,t)+s\mathbf1_{\{t\le q\}}\bigr],
 \qquad \mu\gets\mu+Pv\rho.
\end{equation}
For a merged option use $d_{\rm m}$ instead. All corrections in a
row are constant in $u$. This gives the joint transfer
\[
 \sum_t m_t e_\beta(t)+
 \mathbb E\left[\sum_t m_t d(u,t)+\operatorname{Save}(G)\right]
 \ge\sum_t m_t a_t-\delta w\mu.
\]

Set $\tau=20718779/10^7=2.0718779$. Rational logarithm enclosures
certify that only $a_1$ is below $\tau$:
\begin{equation}\label{eq:first-phase-coefficient-certificate}
 \begin{aligned}
 1.223255846141&<a_1<1.223255846143<\tau,\\
 a_2&>2.07187805487,&a_3&>2.07187805533,\\
 a_4&>2.07187805471,&a_5&>2.07187805502,\\
 a_6&>2.07187805311,&a_7&>2.07187805484.
 \end{aligned}
\end{equation}
Every saving cutoff is at most six, and every signed phase correction
is constant for $t\ge7$. Thus $a_t=e_\beta(t)+\kappa$ for one fixed
$\kappa$ and all $t\ge7$. Since $a_7>\tau$ and $e_\beta$ is
nondecreasing, every later coefficient satisfies $a_t\ge a_7>\tau$.
Define
\begin{equation}\label{eq:first-phase-chi-B}
 \chi=\tau-2\delta(\tau-a_1)-\delta\mu,\qquad
 B=7-\frac{\ln6}{\ln\beta}-\chi.
\end{equation}
These are exactly the coefficients of
Lemma~\ref{lem:capped-rank-credit} with exceptional index $E=\{1\}$.
To invoke that lemma, it remains to retain the zero-route discount
under the same table distribution.

\subsection{Retain the zero-route credit and convert the stage}

Let $D_0(k)$ be the table-averaged length change in
\eqref{eq:first-phase-zero-correction}. The same exact certificate gives
\[
 \begin{array}{c|rrrrrrr}
 k&1&2&3&4&5&6&7\\\hline
 D_0(k)&0&<-.24242&<-.31657&<-.36504&<-.36504&<-.32434&<-.28498
 \end{array}
\]
For all $k\ge7$ the change is $D_0(7)$: regular initial budgets are
at most six, $k_0\le3$, and a merged change depends only on $k_0$
once $k>k_0$. Hence $D_0(k)\le0$ for every possible minimum cover
size. The modified family retains the original joint mean length bound
\[
 \mathbb E[|G|]\le Z_\beta(k_{\min})
 \le\nu_\beta k_{\min}-\frac1{\ln\beta}.
\]
If the zero pattern exists, a supporting witness path of length $r$
gives $k_{\min}\le r$, as in Section~\ref{sec:polynomial-blocks}.
Since $\lambda_\beta>\nu_\beta$, the same restriction argument
retains the doubled zero-route credit $2p_0/\ln\beta$.

The rational certificate also verifies $\chi+2/\ln\beta>\tau$.
Lemma~\ref{lem:capped-rank-credit}, with $\zeta=2/\ln\beta$,
therefore turns the coefficient and intercept bounds into the saving
$\chi b$. This step includes arbitrary zero-pattern mass $p_0$.
Minimizing the rational surrogate is at most averaging it over the
table's schedule distribution. The discrete discovery proof and
\eqref{eq:adaptive-stage} therefore yield
\begin{equation}\label{eq:first-phase-local}
 W_{\rm pat}(S)\le A\Delta+B b+\frac{I-J_*}{\ln2},
 \qquad A=\frac{\beta+1}{\ln\beta}.
\end{equation}
The at-most-two-real-tests degenerate hull construction obeys this
bound too, since $A,B>2$, by the same argument as in the earlier
block proofs.

\begin{theorem}[Improved initial-batch polynomial approximation]
\label{thm:first-phase-expected}
For general supplied binary unit-cost planar line queries and arbitrary
strictly positive rational priors, the initial-batch adaptive algorithm
runs in deterministic polynomial time and satisfies
\[
 c_{\mathrm E}(T(\HH))\le\rho_{\rm first}\OPT_{\mathrm E}(\HH),
 \qquad
 \boxed{\rho_{\rm first}=
 A+\max\{\ln2/h(\delta),B\ln2/\sigma(\delta)\}<5.197569.}
\]
Its running-time exponent has no dependence on the smallest prior.
\end{theorem}
\begin{proof}
Apply Lemma~\ref{lem:stage-conversion} to
\eqref{eq:first-phase-local} and
\eqref{eq:coarse-discovery-cap}. The coarsened transcript is a
function of the full pattern, so $0\le J_*\le I$ as required.
With the parameters in \eqref{eq:first-phase-parameters},
$6<1/\delta<7$ and
\[
 \sigma(\delta)=-6\delta\ln\delta
               -(1-6\delta)\ln(1-6\delta).
\]
The rational atanh bounds \eqref{eq:rational-log-bounds}, after
range reduction of every argument to $[1,2]$, reproduce all the
comparisons above with $N=100$ terms. For completeness, if
$n_t=\lfloor\beta t\rfloor$ and $2\le t\le7$, the values used are
\[
 \begin{aligned}
 e_\beta(t)={}&
 \frac{(\beta+1)(t-1)-(t+n_t)\ln t+\ln(t!)}{\ln\beta}\\
 &-(n_t-3)+\frac{\ln(n_t!/6)}{\ln\beta}.
 \end{aligned}
\]
Together with the eight-row definitions, these give
\[
 \begin{aligned}
 4.379796439664&<\mu<4.379796439666,\\
 1.128579101529&<\chi<1.128579101532,\\
 4.445350596842&<B<4.445350596845,\\
 A+\max\{\ln2/h(\delta),B\ln2/\sigma(\delta)\}
 &<5.197568047723<5.197569.
 \end{aligned}
\]
The reproducible certificate \texttt{next-expected-check.py} checks
these inequalities with exact fractions and outward logarithm bounds;
it also verifies both phase identities and $1{,}950$ finite modified
batch comparisons, including early minimum-cover completion.
The fixed parameter table and its outward interval bounds give a
reproducible certificate for these numerical claims.

The coverage DPs, polynomial finite shift-cell enumeration, constant
number of initial variants, and disjoint-pattern recursion preserve
polynomial time. Logarithms and the table mixture appear only in the
proof. Every actual query remains a supplied endpoint line.
\end{proof}
